%% toc: P2 23 23.1,23.2,23.3,23.4
\section{Computably enumerable sets are \(\Sigma^0_1\)-definable: the arithmetic step of nonaxiomatizability}
\label{sec:p02-w02b-p02-034-ce-definable}

\emph{Purpose.}  Clause~3 of \cref{thm:p02-w02-p02-040-no-effective-calculus}
says that true arithmetic \(\mathrm{TA}\) is not computably enumerable, and
the proof given with that theorem derives its clauses~4 and~5 (no effective
complete calculus for full second-order validity) from clause~3.  The step from computable
enumerability to arithmetic definability is the one place where that
argument leaves logic and enters the theory of machines.  This section
proves that step directly from the Turing-machine definitions of
\cref{sec:turing-machines-01,sec:turing-machines-03}: for each machine an
explicit \(\Sigma^0_1\) formula is written down that holds of \(\bar x\) in
\(\mathcal N\) exactly when the machine halts on the binary numeral of
\(x\).  Clause~3 then follows from \cref{thm:tarski-undefinability}.
Clauses~4 and~5 follow from the same halting formula and an arithmetical
graph for the primitive recursive map \(c\mapsto\ulcorner\mathsf{PA}_2\to\sigma\urcorner\),
again by \cref{thm:tarski-undefinability}; no machine is built and no two
machines are run in sequence.  The handbook's other proof of clause~3
obtains its arithmetic step by citing \cref{thm:p04-ce-sigma-one}, and its
other proof of clauses~4 and~5 builds machines from primitive recursive
functions; \cref{def:p02-w02b-p02-034-boundary} states exactly what those
steps carry and what is proved here without them.

\emph{Inputs.}  Deterministic single-tape Turing machines, configurations,
the one-step relation and computations (\cref{sec:turing-machines-01});
the binary word \(\operatorname{bin}(n)\) and partial computable functions
(\cref{def:partial-computable}); computably enumerable sets as domains of
partial computable functions (\cref{sec:turing-machines-03}); the language
\(\Lang_A\), the structure \(\mathcal N\), numerals, marked bounded
quantifiers and the syntactic classes \(\Sigma^0_n\)
(\cref{sec:decision-hierarchy-02}); the beta formula \(\mathsf B\) of
\cref{lem:goedel-beta}; the substitution lemma \cref{thm:fol-substitution};
for the last result, the set \(\mathrm{TA}\) of
\cref{def:p02-w02-p02-040-so-coding}, the arithmetical truth definitions of
\cref{sec:second-incompleteness-04} and \cref{thm:tarski-undefinability};
for the results on full validity, the primitive recursive functions of
\cref{sec:turing-machines-02}, the graph formulas \(\Gamma_d\) of
\cref{def:p05-pa-pr-graph}, the code of \(\mathsf{PA}_2\to\sigma\) and
clauses~1 and~2 of \cref{thm:p02-w02-p02-040-no-effective-calculus}.
\emph{Conclusions.}
\Cref{thm:p02-w02b-p02-034-step,thm:p02-w02b-p02-034-init,thm:p02-w02b-p02-034-halt,thm:p02-w02b-p02-034-ce-sigma-one,thm:p02-w02b-p02-034-ta-not-ce,thm:p02-w02b-p02-034-pr-graph,thm:p02-w02b-p02-034-reduction,thm:p02-w02b-p02-034-no-calculus}.
\emph{Standing.}  The construction and every proof are the handbook's own
and are written in full in the annex.  Nothing here uses
\cref{thm:p04-ce-sigma-one}, the bounded verification formula
\(\mathsf{Comp}\), the enumeration theorem, the equivalence of Turing
computability with partial recursiveness, the computability of primitive
recursive functions, or the sequential composition of machines.

\begin{proofboundary}{def:p02-w02b-p02-034-boundary}{Dependencies of the nonaxiomatizability argument}
The five clauses of \cref{thm:p02-w02-p02-040-no-effective-calculus}
use two distinct constructions. The syntax recognizers and the map
\(r(c)=\ulcorner\mathsf{PA}_2\to\sigma\urcorner\), for
\(c=\ulcorner\sigma\urcorner\), are primitive recursive.
\Cref{thm:p04-primitive-recursive-total} makes their characteristic
functions and code transformation computable. The strict composition
and minimization clauses of partial recursion, together with
\cref{thm:p04-tm-partial-recursive}, supply the guarded preimage
programs used in clauses~4 and~5. Their exact domains are verified in
\cref{annex:p02-w02-p02-040-no-effective-calculus}.

The arithmetic descriptions of a step and of the input are proved
in \cref{thm:p02-w02b-p02-034-step,thm:p02-w02b-p02-034-init}.
\Cref{thm:p02-w02b-p02-034-halt} combines them into a halting
formula, using the bounded sequence coding of \cref{lem:goedel-beta}.

The machine witnessing partial computability accepts with a valid
output on its domain and diverges outside it, by
\cref{def:partial-computable}. Its halting set is therefore exactly
that domain. \Cref{thm:p02-w02b-p02-034-ce-sigma-one} supplies the
forward enumerability implication used in clause~3. Tarski's theorem
then excludes computable enumerability of true arithmetic.

The formula route of
\cref{thm:p02-w02b-p02-034-pr-graph,thm:p02-w02b-p02-034-reduction,thm:p02-w02b-p02-034-no-calculus}
instead substitutes an arithmetic graph of \(r\) into a halting
formula. It proves clauses~4 and~5 without constructing a composed
machine for that reduction. Both routes concern full second-order
semantics. This section proves the forward enumerability implication;
the converse and the decidability characterization belong to
\cref{thm:p04-ce-sigma-one}.
\end{proofboundary}

\subsection{Numerical codes of configurations}
\label{sec:p02-w02b-p02-034-codes}

\begin{hbconvention}{def:p02-w02b-p02-034-machine}{The fixed machine and its digits}
Throughout the section fix a Turing machine
\(\mathcal T=(Q,\mathcal A_{\mathrm{in}},\mathcal A_{\mathrm{tape}},\delta,q_0,q_{\mathrm{acc}},q_{\mathrm{rej}})\)
in the sense of \cref{sec:turing-machines-01} whose input alphabet contains
the two symbols \(\mathtt 0\) and \(\mathtt 1\).  Typewriter digits are tape
symbols; ordinary digits are numbers.  Let \(m\) be the number of tape
symbols.  Since \(\mathtt 0,\mathtt 1\in\mathcal A_{\mathrm{in}}\) and
\(\sqcup\notin\mathcal A_{\mathrm{in}}\), \(m\geq3\).  Fix
\begin{enumerate}
 \item a bijection \(\operatorname{dg}:\mathcal A_{\mathrm{tape}}\to\{0,\ldots,m-1\}\)
   with \(\operatorname{dg}(\sqcup)=0\), the \emph{digit} of a symbol;
 \item the \emph{base} \(b\defeq m\);
 \item an injection \(\operatorname{st}:Q\to\N\), the \emph{state number}.
\end{enumerate}
The halting states are \(q_{\mathrm{acc}}\) and \(q_{\mathrm{rej}}\); the
initial state \(q_0\) may be one of them.
\end{hbconvention}

\begin{hbdefinition}{def:p02-w02b-p02-034-code}{Left code, right code and code of a configuration}
For a configuration \(C=(q,T,h)\) of \(\mathcal T\) put
\[
 \operatorname{left}(C)\defeq\sum_{j\in\N}\operatorname{dg}(T(h-1-j))\,b^j,
 \qquad
 \operatorname{right}(C)\defeq\sum_{j\in\N}\operatorname{dg}(T(h+j))\,b^j,
\]
and \(\operatorname{enc}(C)\defeq(\operatorname{st}(q),\operatorname{left}(C),\operatorname{right}(C))\in\N^3\).
Both sums are finite, because \(T\) has finite nonblank support and
\(\operatorname{dg}(\sqcup)=0\).  The base-\(b\) digits of
\(\operatorname{right}(C)\), least significant first, are the digits of the
scanned cell and the cells to its right; those of \(\operatorname{left}(C)\)
are the digits of the cells to the left of the head, nearest first.  The
code forgets the head position \(h\): configurations that differ by a
translation of the tape together with the head have the same code.
\end{hbdefinition}

\begin{hbdefinition}{def:p02-w02b-p02-034-formulas}{The step, input and halting formulas}
Write \(\mathsf B(w,i,v)\) for the bounded formula of \cref{lem:goedel-beta},
whose bound variables are renamed apart from every variable displayed below,
and \(\beta(w,i)\) for the unique \(v\) it relates to \(w,i\) in
\(\mathcal N\).  For a term \(\tau\) in place of \(i\) or \(v\),
\(\mathsf B(w,\tau,v)\) and \(\mathsf B(w,i,\tau)\) are the substitution
instances.

\emph{Step.}  Let \(D\) be the finite set of pairs \((q,a)\) with
\(q\in Q\setminus\{q_{\mathrm{acc}},q_{\mathrm{rej}}\}\) and
\(a\in\mathcal A_{\mathrm{tape}}\).  For \((q,a)\in D\) with
\(\delta(q,a)=(q',c,d)\), let \(\Delta_{q,a}\) be, if \(d=R\),
\[
\begin{aligned}
 &p=\overline{\operatorname{st}(q)}\land r=(\bar b\times u)+\overline{\operatorname{dg}(a)}
   \land p'=\overline{\operatorname{st}(q')}\\
 &\quad\land l'=(\bar b\times l)+\overline{\operatorname{dg}(c)}
   \land r'=u,
\end{aligned}
\]
and, if \(d=L\),
\[
\begin{aligned}
 &p=\overline{\operatorname{st}(q)}\land r=(\bar b\times u)+\overline{\operatorname{dg}(a)}
   \land p'=\overline{\operatorname{st}(q')}\\
 &\quad\land\bigvee_{g<b}\bigl(l=(\bar b\times l')+\bar g\land
   r'=\bar g+\bar b\times(\overline{\operatorname{dg}(c)}+(\bar b\times u))\bigr),
\end{aligned}
\]
and put
\(\mathsf{Step}_{\mathcal T}(p,l,r,p',l',r')\defeq\exists u\bigvee_{(q,a)\in D}\Delta_{q,a}\).
An empty disjunction is \(\bottom\).

\emph{Input.}  \(\mathsf{Init}_{\mathcal T}(x,r)\) is
\[
\begin{aligned}
 \exists k\,\exists s_1\,\exists s_2\,\exists s_3\,\exists s_4\,\Bigl[\,
 &\mathsf B(s_2,\bar0,\bar0)\land\mathsf B(s_3,\bar0,\bar0)\land\mathsf B(s_4,\bar0,\bar1)\\
 &\land\forall j<Sk\,\exists e\,\exists v\,\exists v'\,\exists t\,\exists t'\,\exists\pi\,\exists\pi'\,
 \bigl(\mathsf B(s_1,j,e)\\
 &\qquad\land\mathsf B(s_2,j,v)\land\mathsf B(s_2,Sj,v')\land
   \mathsf B(s_3,j,t)\land\mathsf B(s_3,Sj,t')\\
 &\qquad\land\mathsf B(s_4,j,\pi)\land\mathsf B(s_4,Sj,\pi')\land\pi'=\bar b\times\pi\\
 &\qquad\land\bigl((e=\bar0\land v'=v+v\land t'=t+(\overline{\operatorname{dg}(\mathtt 0)}\times\pi))\\
 &\qquad\quad\lor(e=\bar1\land v'=S(v+v)\land t'=t+(\overline{\operatorname{dg}(\mathtt 1)}\times\pi))\bigr)\bigr)\\
 &\land\mathsf B(s_2,Sk,x)\land\mathsf B(s_3,Sk,r)\land\bigl(k=\bar0\lor\mathsf B(s_1,\bar0,\bar1)\bigr)\Bigr].
\end{aligned}
\]
\emph{Halting.}  \(\mathsf{Halt}_{\mathcal T}(x)\) is
\[
\begin{aligned}
 \exists n\,\exists w_1\,\exists w_2\,\exists w_3\,\exists r_0\,\Bigl[\,
 &\mathsf{Init}_{\mathcal T}(x,r_0)\land\mathsf B(w_1,\bar0,\overline{\operatorname{st}(q_0)})\\
 &\land\mathsf B(w_2,\bar0,\bar0)\land\mathsf B(w_3,\bar0,r_0)\\
 &\land\forall k<n\,\exists p\,\exists l\,\exists r\,\exists p'\,\exists l'\,\exists r'\,
   \bigl(\mathsf B(w_1,k,p)\\
 &\qquad\land\mathsf B(w_2,k,l)\land\mathsf B(w_3,k,r)\\
 &\qquad\land\mathsf B(w_1,Sk,p')\land\mathsf B(w_2,Sk,l')\land\mathsf B(w_3,Sk,r')\\
 &\qquad\land\mathsf{Step}_{\mathcal T}(p,l,r,p',l',r')\bigr)\\
 &\land\exists p_\ast\,\bigl(\mathsf B(w_1,n,p_\ast)
   \land(p_\ast=\overline{\operatorname{st}(q_{\mathrm{acc}})}\\
 &\qquad\lor p_\ast=\overline{\operatorname{st}(q_{\mathrm{rej}})})\bigr)\Bigr].
\end{aligned}
\]
All variables shown are distinct, and the bound variables of
\(\mathsf{Init}_{\mathcal T}\) and \(\mathsf{Step}_{\mathcal T}\) are renamed
apart from those of \(\mathsf{Halt}_{\mathcal T}\).  The only free variable
of \(\mathsf{Halt}_{\mathcal T}\) is \(x\).  The finite disjunctions over
\(D\) and over \(g<b\) are metalinguistic expansions fixed by
\(\mathcal T\) and the choices of \cref{def:p02-w02b-p02-034-machine}.
\index{computation, arithmetical description of}
\end{hbdefinition}

The three formulas transcribe three facts.  A single step changes only the
lowest digits of the two codes (\cref{thm:p02-w02b-p02-034-step}).  The
initial right code is obtained from \(x\) by reading \(\operatorname{bin}(x)\)
from its most significant bit, accumulating the binary value of the prefix
read and the base-\(b\) code of the cells written so far
(\cref{thm:p02-w02b-p02-034-init}).  A halting computation is a finite
sequence of codes, which the beta function stores in three numbers
\(w_1,w_2,w_3\) (\cref{thm:p02-w02b-p02-034-halt}).

\subsection{Results}
\label{sec:p02-w02b-p02-034-results}

\begin{hblemma}{thm:p02-w02b-p02-034-step}{proved=annex:p02-w02b-p02-034-step}{One step of the machine as arithmetic}
Let \(C\) be a configuration of \(\mathcal T\) and \(p',l',r'\in\N\).
\begin{enumerate}
 \item If \(C\longrightarrow_{\mathcal T}C'\), then
   \(\mathcal N\entails\mathsf{Step}_{\mathcal T}(\overline{\operatorname{enc}(C)},\overline{\operatorname{enc}(C')})\).
 \item If
   \(\mathcal N\entails\mathsf{Step}_{\mathcal T}(\overline{\operatorname{enc}(C)},\bar p',\bar l',\bar r')\),
   then \(C\) is not halting, it has a successor \(C'\), and
   \((p',l',r')=\operatorname{enc}(C')\).
\end{enumerate}
Here \(\overline{\operatorname{enc}(C)}\) stands for the three numerals of
the three entries.
\end{hblemma}

\begin{hblemma}{thm:p02-w02b-p02-034-init}{proved=annex:p02-w02b-p02-034-init}{The initial configuration as arithmetic}
For \(x\in\N\) let \(C^x_0=(q_0,T_{\operatorname{bin}(x)},0)\) be the
initial configuration of \(\mathcal T\) on \(\operatorname{bin}(x)\).  Then
\(\operatorname{left}(C^x_0)=0\), and for every \(r\in\N\),
\(\mathcal N\entails\mathsf{Init}_{\mathcal T}(\bar x,\bar r)\) iff
\(r=\operatorname{right}(C^x_0)\).
\end{hblemma}

\begin{hbtheorem}{thm:p02-w02b-p02-034-halt}{proved=annex:p02-w02b-p02-034-halt}{Halting is \(\Sigma^0_1\)-definable}
\(\mathsf{Halt}_{\mathcal T}(x)\) is a \(\Sigma^0_1\) formula in the sense of
\cref{sec:decision-hierarchy-02}, and for every \(x\in\N\),
\[
 \mathcal N\entails\mathsf{Halt}_{\mathcal T}(\bar x)
 \quad\Longleftrightarrow\quad
 \text{the computation of }\mathcal T\text{ on }\operatorname{bin}(x)\text{ halts}.
\]
\end{hbtheorem}

\begin{hbcorollary}{thm:p02-w02b-p02-034-ce-sigma-one}{proved=annex:p02-w02b-p02-034-ce-sigma-one}{Every computably enumerable set is \(\Sigma^0_1\)-definable}
If \(A\subseteq\N\) is computably enumerable, then
\(A\in\boldsymbol\Sigma^0_1\): there is a \(\Sigma^0_1\) formula
\(\varphi(x)\) with \(a\in A\) iff \(\mathcal N\entails\varphi(\bar a)\) for
every \(a\in\N\).  If \(A\) is the domain of the partial function computed
by \(\mathcal T\), then \(\varphi\) can be taken to be
\(\mathsf{Halt}_{\mathcal T}\).
\end{hbcorollary}

This is the direction of \cref{thm:p04-ce-sigma-one} from computable
enumerability to \(\boldsymbol\Sigma^0_1\), proved here by a route that does
not pass through the enumeration theorem.  The corollary
above is a full proof of that direction which does not depend on how
\cref{thm:p04-ce-sigma-one} itself is proved
(\cref{def:p02-w02b-p02-034-boundary}).  The converse direction and the clause on
\(\boldsymbol\Delta^0_1\) are not proved in this section and are not used by
it.

\begin{hbcorollary}{thm:p02-w02b-p02-034-ta-not-ce}{proved=annex:p02-w02b-p02-034-ta-not-ce}{True arithmetic is not computably enumerable}
The set \(\mathrm{TA}\) of codes of first-order \(\Lang_A\)-sentences true in
\(\mathcal N\) (\cref{def:p02-w02-p02-040-so-coding}) is not computably
enumerable.  This is clause~3 of
\cref{thm:p02-w02-p02-040-no-effective-calculus}.
\end{hbcorollary}

The preceding corollary combines the halting-formula construction
with Tarski's theorem. This is also the arithmetic argument in
clause~3 of \cref{annex:p02-w02-p02-040-no-effective-calculus}.
No enumeration-theorem premise is needed.

The proofs of clauses~4 and~5 in
\cref{annex:p02-w02-p02-040-no-effective-calculus} use a guarded
partial-recursive composition, realized by
\cref{thm:p04-tm-partial-recursive}. The three results below give a
second route, replacing that composed computation by a formula.  A computably enumerable set of codes is
defined in \(\mathcal N\) by a halting formula
(\cref{thm:p02-w02b-p02-034-ce-sigma-one}); the primitive recursive map
\(r\) of clause~1 has an arithmetical graph
(\cref{thm:p02-w02b-p02-034-pr-graph}); substituting the second into the
first defines, in arithmetic, the set of first-order sentences \(\sigma\)
for which \(\mathsf{PA}_2\to\sigma\) lies in the given set.  When the given
set agrees with full validity on these implications, clause~2 makes that
definition an arithmetical truth definition
(\cref{thm:p02-w02b-p02-034-reduction}).

\begin{hblemma}{thm:p02-w02b-p02-034-pr-graph}{proved=annex:p02-w02b-p02-034-pr-graph}{Graphs of primitive recursive functions}
Let \(f:\N^k\to\N\) be primitive recursive in the sense of
\cref{sec:turing-machines-02}.  Then \(f=f_d\) for some primitive recursive
description \(d\) of arity \(k\) in the sense of \cref{def:p05-pa-pr-graph}.
The graph formula \(\Gamma_d(x_1,\ldots,x_k,y)\) of that definition
satisfies, for all \(a_1,\ldots,a_k,b\in\N\),
\[
 \mathcal N\entails\Gamma_d(\bar a_1,\ldots,\bar a_k,\bar b)
 \quad\Longleftrightarrow\quad
 f(a_1,\ldots,a_k)=b.
\]
\end{hblemma}

The displayed equivalence is the standard-model clause of
\cref{thm:p05-pa-pr-graphs}; the annex writes out, step by step, the
induction on the description that the proof of that theorem records in one
paragraph.  Nothing about derivability in \(\PA\) or \(\mathsf Q\) is used.

\begin{hblemma}{thm:p02-w02b-p02-034-reduction}{proved=annex:p02-w02b-p02-034-reduction}{No arithmetically definable set agrees with full validity on the implications from \(\mathsf{PA}_2\)}
There is no set \(S\subseteq\N\) with both of the following properties:
\begin{enumerate}
 \item some \(\Lang_A\)-formula \(\varphi(x)\) whose only free variable is
   \(x\) satisfies, for every \(c\in\N\), \(c\in S\) iff
   \(\mathcal N\entails\varphi(\bar c)\);
 \item for every first-order \(\Lang_A\)-sentence \(\sigma\),
   \(\ulcorner\mathsf{PA}_2\to\sigma\urcorner\in S\) iff
   \(\entails_{\mathrm{full}}(\mathsf{PA}_2\to\sigma)\).
\end{enumerate}
\end{hblemma}

\begin{hbcorollary}{thm:p02-w02b-p02-034-no-calculus}{proved=annex:p02-w02b-p02-034-no-calculus}{Full validity is not computably enumerable}
\begin{enumerate}
 \item \(\mathrm{Val}_2\) is not computably enumerable.
 \item There is no calculus \(\mathsf D\) whose set of theorem codes, in the
   coding of \cref{def:p02-w02-p02-040-so-coding}, is computably enumerable
   and which satisfies
   \(\mathsf D\derives\sigma\Leftrightarrow\entails_{\mathrm{full}}\sigma\)
   for every second-order sentence \(\sigma\) over \(\{0,S,+,\times\}\).
\end{enumerate}
These are clauses~4 and~5 of
\cref{thm:p02-w02-p02-040-no-effective-calculus}.
\end{hbcorollary}

The chain for clause~4 reads: a computably enumerable \(\mathrm{Val}_2\) is
defined in \(\mathcal N\) by \(\mathsf{Halt}_{\mathcal T}\) for a suitable
\(\mathcal T\); the formula
\(\exists y\,(\Gamma_r(x,y)\land\mathsf{Halt}_{\mathcal T}(y))\), with
\(\Gamma_r\) a graph formula for \(r\), then holds of the code of a
first-order sentence \(\sigma\) exactly when
\(\entails_{\mathrm{full}}(\mathsf{PA}_2\to\sigma)\), that is, by clause~2,
exactly when \(\mathcal N\entails\sigma\); and so it is an arithmetical truth
definition.  For clause~5 the theorem codes of \(\mathsf D\) take the place
of \(\mathrm{Val}_2\).  The decidability of the sentence test, which the
machine proof needs, plays no part: a truth definition is tested only on
codes of sentences.

\begin{hbexample}{def:p02-w02b-p02-034-example}{A one-step machine traced in arithmetic}
Let \(\mathcal T\) have states \(q_0,q_{\mathrm{acc}},q_{\mathrm{rej}}\), input
alphabet \(\{\mathtt 0,\mathtt 1\}\), tape alphabet
\(\{\sqcup,\mathtt 0,\mathtt 1\}\), and \(\delta(q_0,a)=(q_{\mathrm{acc}},a,R)\)
for each tape symbol \(a\).  It halts on every input after one step.  Take
\(\operatorname{dg}(\sqcup)=0\), \(\operatorname{dg}(\mathtt 0)=1\),
\(\operatorname{dg}(\mathtt 1)=2\), so \(b=3\), and
\(\operatorname{st}(q_0)=0\), \(\operatorname{st}(q_{\mathrm{acc}})=1\),
\(\operatorname{st}(q_{\mathrm{rej}})=2\).

Input \(x=2\), so \(\operatorname{bin}(2)=\mathtt{10}\) occupies cells \(0\)
and \(1\).  In \(\mathsf{Init}_{\mathcal T}\) take \(k=1\) and the sequences
\(e=(1,0)\), \(v=(0,1,2)\), \(\pi=(1,3,9)\) and
\(t=(0,0+2\cdot1,2+1\cdot3)=(0,2,5)\).  Every clause holds: \(e_0=1\), the
binary values end at \(v_2=2=x\), and \(t_2=5\).  Directly,
\(\operatorname{right}(C^2_0)=\operatorname{dg}(\mathtt 1)+\operatorname{dg}(\mathtt 0)\cdot3=5\)
and \(\operatorname{left}(C^2_0)=0\), so \(\operatorname{enc}(C^2_0)=(0,0,5)\).

One step: \(5=3\cdot1+2\), so \(u=1\) and the scanned digit is
\(2=\operatorname{dg}(\mathtt 1)\).  The disjunct \(\Delta_{q_0,\mathtt 1}\) of
\(\mathsf{Step}_{\mathcal T}\), with \(\delta(q_0,\mathtt 1)=(q_{\mathrm{acc}},\mathtt 1,R)\),
gives \(p'=1\), \(l'=3\cdot0+2=2\) and \(r'=1\).  On the tape, the head has
moved to cell \(1\), which holds \(\mathtt 0\) (right code \(1\)), and cell \(0\)
to its left holds \(\mathtt 1\) (left code \(2\)): the code of the successor is
\((1,2,1)\), as \cref{thm:p02-w02b-p02-034-step} predicts.  With \(n=1\) and
\(w_1,w_2,w_3\) coding \((0,1)\), \((0,2)\), \((5,1)\), the final clause of
\(\mathsf{Halt}_{\mathcal T}\) holds because \(1=\operatorname{st}(q_{\mathrm{acc}})\).
So \(\mathcal N\entails\mathsf{Halt}_{\mathcal T}(\bar2)\).
\end{hbexample}

\begin{hbwarning}{def:p02-w02b-p02-034-limits}{What the construction gives and what it does not}
\begin{enumerate}
 \item The equivalence of \cref{thm:p02-w02b-p02-034-halt} is truth in the
   standard structure \(\mathcal N\).  Nothing is claimed about derivability of
   \(\mathsf{Halt}_{\mathcal T}(\bar x)\) or its negation in \(\mathsf Q\) or in
   \(\PA\); the representability results used for incompleteness are a
   separate matter (\cref{sec:p05-representability}).
 \item This section proves the direction from computable enumerability
   to \(\boldsymbol\Sigma^0_1\). The converse and the
   \(\boldsymbol\Delta^0_1\) characterization are the separate
   assertions of \cref{thm:p04-ce-sigma-one}; neither is used in
   the present nonaxiomatizability argument.
 \item The code of a configuration forgets the head position, and its digits
   do not record where the nonblank part of the tape ends.  Neither matters:
   the transition depends only on the state and the scanned symbol, and
   blank cells have digit \(0\).  A coding that gave the blank a nonzero digit
   would make the sums of \cref{def:p02-w02b-p02-034-code} infinite.
 \item The argument concerns full semantics.  The calculus of
   \cref{def:p02-sol-derived-calculus} is complete for Henkin semantics
   (\cref{thm:p02-sol-derived-adequacy}); clause~3 and the resulting
   nonaxiomatizability say nothing against that completeness, which is for a
   different consequence relation.
 \item The formula proof of
   \cref{thm:p02-w02b-p02-034-no-calculus} needs only the arithmetic
   graph of the reduction, not a machine deciding sentencehood.
   The latter decision procedure, used in the computational route,
   follows from the explicit primitive-recursive syntax tests and
   \cref{thm:p04-primitive-recursive-total}.
\end{enumerate}
\end{hbwarning}

\emph{Relationships.}  \Cref{thm:p02-w02b-p02-034-ce-sigma-one} is the
direction of \cref{thm:p04-ce-sigma-one} used by the
corollary to \cref{thm:tarski-undefinability} and by clause~3 of
\cref{thm:p02-w02-p02-040-no-effective-calculus}, proved without the
enumeration theorem or \(\mathsf{Comp}\)
(\cref{def:p02-w02b-p02-034-boundary}).  The proof of clauses~4
and~5 given in \cref{sec:p02-w02-p02-041-direct} avoids true arithmetic
altogether but passes from Turing machines to partial recursive functions
through \cref{thm:p04-tm-partial-recursive}; the present route instead stays
with Turing machines and needs no such passage.  For clauses~4 and~5 it also
builds no machine from primitive recursive functions: the primitive
recursive reduction \(r\) enters only through its arithmetical graph
(\cref{thm:p02-w02b-p02-034-pr-graph,thm:p02-w02b-p02-034-no-calculus}).  A much stronger
definability statement holds: a set of positive integers is Diophantine iff
it is recursively enumerable \cite[\S8, Theorem~8.1, p.~264]{Davis1973},
where recursive enumerability is defined there through recursive functions
\cite[\S8, Definition~8.1, p.~264]{Davis1973}, so that a single existential
block over a polynomial equation suffices, with no bounded quantifier.  That
theorem is not used here.

\section{Two conventions of Church's formulation: autonymy and proof by induction}
\label{sec:p02-w02b-p02-034-church-conventions}

\emph{Purpose.}  \Cref{sec:p02-w02-p02-040-church} presents the system of
Church's paper of 1940 \cite{Church1940}: its symbols, formation rules,
abbreviations, rules, axioms and selected derivations.  Two conventions of
the paper govern how its statements and proofs are to be read, and they are
fixed here.  The first concerns the language in which the system is
described; the second names a way of presenting proofs that rests on the
induction theorem VIII (\cref{thm:p02-w02-p02-040-church-induction}).
\emph{Inputs.}  \Cref{def:p02-w02-p02-040-church-wff,def:p02-w02-p02-040-church-abbrev,def:p02-w02-p02-040-church-rules},
the deduction theorem \cref{thm:p02-w02-p02-040-church-deduction}, theorems
13\(^\alpha\) and 16\(^\alpha\) of \cref{thm:p02-w02-p02-040-church-12-21},
theorem 24\(^\alpha\) and the induction theorem VIII of
\cref{thm:p02-w02-p02-040-church-induction}, and the propositional input (P)
of \cref{def:p02-w02-p02-040-church-prop}, which
\cref{thm:p02-w02-p02-040-church-prop-complete} proves.
\emph{Conclusions.}  \Cref{def:p02-w02b-p02-034-autonymy,def:p02-w02b-p02-034-induction-proof}
and the worked derivation of \cref{def:p02-w02b-p02-034-induction-example}.
Notation follows the source and \cref{sec:p02-w02-p02-040-church}.

\begin{hbconvention}{def:p02-w02b-p02-034-autonymy}{Formal symbols as names of themselves}
Church describes the formal system in English, which serves as the syntax
language.  In statements about the system the source uses bold capital
letters as variables ranging over well-formed formulas and bold small
letters as variables ranging over variables of the system, with a subscript
giving the type.  Beyond these, it adopts the usage, which it calls
customary and self-explanatory, under which every symbol of the formal
language, written in the syntax language, is a name of that symbol itself,
and writing expressions next to one another denotes the expression obtained
by writing what they denote one after another
\cite[\S2, pp.~57--58]{Church1940}.

So a display such as \((\lambda\mathbf x_\beta\mathbf M_\alpha)\) in the
second formation rule (\cref{def:p02-w02-p02-040-church-wff}) is a name, in
the syntax language, of a formula: the formula consisting of the symbol
\((\), then the symbol \(\lambda\), then the variable that
\(\mathbf x_\beta\) denotes, then the symbols of the formula that
\(\mathbf M_\alpha\) denotes in order, then the symbol \()\).  The
displayed symbols \((\), \(\lambda\) and \()\) name themselves; the bold
letters are syntactic variables and do not.  A display that uses the
abbreviations of \cref{def:p02-w02-p02-040-church-abbrev} names the formula
obtained by eliminating them.
\end{hbconvention}

\begin{hbexample}{def:p02-w02b-p02-034-autonymy-example}{Reading a formation rule autonymously}
Take \(\mathbf x_\beta\) to be the variable \(a_\iota\) and
\(\mathbf M_\alpha\) to be the one-symbol formula \(a_\iota\), a
well-formed formula of type \(\iota\) by the first formation rule.  Under the
convention, \((\lambda\mathbf x_\beta\mathbf M_\alpha)\) then denotes the
formula whose primitive symbols are, in order, \((\), \(\lambda\),
\(a_\iota\), \(a_\iota\), \()\): five symbols, written
\((\lambda a_\iota a_\iota)\).  The second formation rule makes it a
well-formed formula of type \(\iota\iota\).  In this last display every
symbol names itself, because no syntactic variable occurs in it; the
occurrences of \(a_\iota\) are the variable of the system, not a variable of
the syntax language.  By the abbreviation \(I_{\alpha\alpha}\) of
\cref{def:p02-w02-p02-040-church-abbrev}, which is \(\lambda x_\alpha x_\alpha\),
\(I_{\iota\iota}\) names \((\lambda x_\iota x_\iota)\), a different formula
that differs from \((\lambda a_\iota a_\iota)\) only by Rule~I.
\end{hbexample}

Two limits follow from the convention as stated.  It applies only where
syntactic variables and formal symbols are typographically distinct, which
the source secures by reserving bold type for the syntactic variables.  And
it concerns names of formulas within the syntax language; it is different
from the G\"odel numbering of \cref{sec:arithmetic-coding-04}, under which
\(\ulcorner\sigma\urcorner\) is a number assigned to a formula, and from the
numeral of that number, which is a term of the object language.

\begin{hbdefinition}{def:p02-w02b-p02-034-induction-proof}{Proof by induction on a variable}
Let \(\mathbf x_{\alpha'}\) be a variable of type \(\alpha'\).  A
\emph{proof by induction on the variable} \(\mathbf x_{\alpha'}\) (also
called a proof by mathematical, or by complete, induction on it) is a proof
that is abbreviated, or can be abbreviated, by using the induction theorem
VIII as a derived rule \cite[\S6, p.~64]{Church1940}: from derivations of
\(\mathbf F_{o\alpha'}0_{\alpha'}\) on assumptions
\(\mathbf A^1_o,\ldots,\mathbf A^n_o\) and of
\(\mathbf F_{o\alpha'}(S_{\alpha'\alpha'}\mathbf x_{\alpha'})\) on the
assumptions
\[
 \mathbf A^1_o,\ldots,\mathbf A^n_o,\ N_{o\alpha'}\mathbf x_{\alpha'},\ \mathbf F_{o\alpha'}\mathbf x_{\alpha'},
\]
it passes to
\[
 \mathbf A^1_o,\ldots,\mathbf A^n_o\derives N_{o\alpha'}\mathbf x_{\alpha'}\supset\mathbf F_{o\alpha'}\mathbf x_{\alpha'}.
\]
The side condition of VIII is part of the rule: \(\mathbf x_{\alpha'}\) is
not free in any \(\mathbf A^i_o\) nor in \(\mathbf F_{o\alpha'}\).
\end{hbdefinition}

VIII is a syntactical theorem, a statement about the system proved in the
syntax language, and not a rule of inference of the system
(\cref{def:p02-w02-p02-040-church-rules} lists six rules, and VIII is none of
them).  A proof by induction is therefore not itself a formal proof: it
establishes that a formal proof exists.  The source obtains VIII from
theorem 24\(^\alpha\) and the deduction theorem \cite[\S6, p.~64]{Church1940},
and the proof of VIII in \cref{annex:p02-w02-p02-040-church-induction} is
the recipe for removing the abbreviation: apply the deduction theorem to the
derivation of the step, generalize upon \(\mathbf x_{\alpha'}\) by VI\('\),
substitute \(\mathbf F_{o\alpha'}\) for \(f_{o\alpha'}\) in 24\(^\alpha\), and
detach twice by Rule~V.  The two side conditions are exactly what make the
generalization and the substitution legitimate.

\begin{hbexample}{def:p02-w02b-p02-034-induction-example}{A proof by induction: every natural number is zero or a successor}
Fix a type \(\alpha\) and choose distinct variables \(m_{\alpha'}\) and
\(k_{\alpha'}\) that occur in none of the expansions of \(0_{\alpha'}\),
\(S_{\alpha'\alpha'}\), \(N_{o\alpha'}\), \(=\) and \(\exists\) in
\cref{def:p02-w02-p02-040-church-abbrev}, and in none of the formulas
13\(^{\alpha'}\) and 16\(^{\alpha'}\).  Let
\[
 \mathbf C_o\;=\;[m_{\alpha'}=0_{\alpha'}]\vee(\exists k_{\alpha'})[N_{o\alpha'}k_{\alpha'}\,[m_{\alpha'}=S_{\alpha'\alpha'}k_{\alpha'}]],
 \qquad
 \mathbf F_{o\alpha'}\;=\;\lambda m_{\alpha'}\mathbf C_o,
\]
the juxtaposition inside the outer square brackets of the existential being
conjunction.  The claim is \(\derives N_{o\alpha'}m_{\alpha'}\supset\mathbf C_o\),
by induction on \(m_{\alpha'}\) with no assumptions (\(n=0\)).

\emph{Side conditions.}  \(m_{\alpha'}\) is bound in \(\mathbf F_{o\alpha'}\),
so it is not free there, and there are no assumptions.  The bound variables
of \(\mathbf C_o\) are \(k_{\alpha'}\) and those of the expansions, and none
of them is \(m_{\alpha'}\); this is the proviso of Rules~II and~III used
below.

\emph{Base.}  Substituting \(0_{\alpha'}\), which has no free variables, for
\(x_{\alpha'}\) in 16\(^{\alpha'}\) by IV\('\) gives
\(\derives0_{\alpha'}=0_{\alpha'}\).  An instance of the tautology
\(p\supset p\vee q\) by (P), and Rule~V, give
\(\derives[0_{\alpha'}=0_{\alpha'}]\vee(\exists k_{\alpha'})[N_{o\alpha'}k_{\alpha'}\,[0_{\alpha'}=S_{\alpha'\alpha'}k_{\alpha'}]]\).
This formula is the result of substituting \(0_{\alpha'}\) for
\(m_{\alpha'}\) throughout \(\mathbf C_o\), so Rule~III gives
\(\derives\mathbf F_{o\alpha'}0_{\alpha'}\).

\emph{Step.}  The assumptions are \(N_{o\alpha'}m_{\alpha'}\) and
\(\mathbf F_{o\alpha'}m_{\alpha'}\); the second is not needed.
\begin{enumerate}
 \item By IV\('\) on 16\(^{\alpha'}\),
   \(\derives S_{\alpha'\alpha'}m_{\alpha'}=S_{\alpha'\alpha'}m_{\alpha'}\).
   With the assumption \(N_{o\alpha'}m_{\alpha'}\), an instance of
   \(p\supset.\,q\supset pq\) by (P) and Rule~V twice give
   \(N_{o\alpha'}m_{\alpha'}\derives[N_{o\alpha'}m_{\alpha'}\,[S_{\alpha'\alpha'}m_{\alpha'}=S_{\alpha'\alpha'}m_{\alpha'}]]\).
 \item In 13\(^{\alpha'}\), Rule~I changes the bound variable
   \(x_{\alpha'}\) of the existential to \(k_{\alpha'}\), which does not occur
   there.  IV\('\) then substitutes \(m_{\alpha'}\) for \(y_{\alpha'}\) and
   \(\mathbf G=\lambda k_{\alpha'}[N_{o\alpha'}k_{\alpha'}\,[S_{\alpha'\alpha'}m_{\alpha'}=S_{\alpha'\alpha'}k_{\alpha'}]]\)
   for \(f_{o\alpha'}\); the only free variable of \(\mathbf G\) is
   \(m_{\alpha'}\), which is not bound in the formula.  Rule~II reduces
   the parts \(\mathbf G\,m_{\alpha'}\) and \(\mathbf G\,k_{\alpha'}\), no
   bound variable of the body of \(\mathbf G\) being \(k_{\alpha'}\) or
   \(m_{\alpha'}\).  The result is the theorem
   \([N_{o\alpha'}m_{\alpha'}\,[S_{\alpha'\alpha'}m_{\alpha'}=S_{\alpha'\alpha'}m_{\alpha'}]]\supset
   (\exists k_{\alpha'})[N_{o\alpha'}k_{\alpha'}\,[S_{\alpha'\alpha'}m_{\alpha'}=S_{\alpha'\alpha'}k_{\alpha'}]]\).
 \item Rule~V with items~1 and~2, then an instance of
   \(q\supset p\vee q\) by (P) and Rule~V, give
   \(N_{o\alpha'}m_{\alpha'}\derives[S_{\alpha'\alpha'}m_{\alpha'}=0_{\alpha'}]\vee(\exists k_{\alpha'})[N_{o\alpha'}k_{\alpha'}\,[S_{\alpha'\alpha'}m_{\alpha'}=S_{\alpha'\alpha'}k_{\alpha'}]]\).
   This is the result of substituting \(S_{\alpha'\alpha'}m_{\alpha'}\) for
   \(m_{\alpha'}\) throughout \(\mathbf C_o\); the only free variable of
   \(S_{\alpha'\alpha'}m_{\alpha'}\) is \(m_{\alpha'}\), so Rule~III gives
   \(N_{o\alpha'}m_{\alpha'},\mathbf F_{o\alpha'}m_{\alpha'}\derives\mathbf F_{o\alpha'}(S_{\alpha'\alpha'}m_{\alpha'})\).
\end{enumerate}

\emph{Conclusion.}  VIII gives
\(\derives N_{o\alpha'}m_{\alpha'}\supset\mathbf F_{o\alpha'}m_{\alpha'}\), and
Rule~II applied to the part \(\mathbf F_{o\alpha'}m_{\alpha'}\) gives
\(\derives N_{o\alpha'}m_{\alpha'}\supset\mathbf C_o\).  Only Axioms
1--6\(^\alpha\), at the types that occur, enter, through (P),
13\(^{\alpha'}\), 16\(^{\alpha'}\) and, inside VIII, 24\(^\alpha\).  Removing the abbreviation
by the recipe above yields a formal proof whose last steps are
24\(^\alpha\) with \(\mathbf F_{o\alpha'}\) for \(f_{o\alpha'}\), followed by
two applications of Rule~V.
\end{hbexample}

\begin{hbwarning}{def:p02-w02b-p02-034-church-open}{What these conventions leave unproved}
Neither convention adds a result to the system.  Two claims that the source
reports without proof keep that standing here: that the same-type
predecessor \(P_{\alpha'\alpha'}\) cannot be defined without descriptions
(\cref{def:p02-w02-p02-041-church-predecessor}), and the two claims of the
source's footnote~12 that Axiom~8 can be replaced by 30\(^\iota\) together
with 30\(^{\iota'}\) in the presence of Axioms 1--6\(^\alpha\), or by
30\(^\iota\) alone in the presence of Axioms 1--7 and 9\(^\alpha\)
\cite[\S6, p.~66, n.~12]{Church1940}.  The handbook proves none of these.
\end{hbwarning}
